\documentclass[12pt]{article}

% Idioma y codificación
\usepackage[spanish, es-tabla]{babel}       %es-tabla para que se titule "Tabla"
\usepackage[utf8]{inputenc}

% Márgenes
\usepackage[a4paper,top=3cm,bottom=2.5cm,left=3cm,right=3cm]{geometry}

% Comentarios de bloque
\usepackage{verbatim}

% Paquetes de links
\usepackage[hidelinks]{hyperref}    % Permite enlaces
\usepackage{url}                    % redirecciona a la web

% Más opciones para enumeraciones
\usepackage{enumitem}

% Personalizar la portada
\usepackage{titling}


% Paquetes de tablas
\usepackage{multirow}


%------------------------------------------------------------------------

%Paquetes de figuras
\usepackage{caption}
\usepackage{subcaption} % Figuras al lado de otras
\usepackage{float}      % Poner figuras en el sitio indicado H.


% Paquetes de imágenes
\usepackage{graphicx}       % Paquete para añadir imágenes
\usepackage{transparent}    % Para manejar la opacidad de las figuras

% Paquete para usar colores
\usepackage[dvipsnames]{xcolor}
\usepackage{pagecolor}      % Para cambiar el color de la página

% Habilita tamaños de fuente mayores
\usepackage{fix-cm}


%------------------------------------------------------------------------

% Paquetes de matemáticas
\usepackage{mathtools, amsfonts, amssymb, mathrsfs}
\usepackage[makeroom]{cancel}     % Simplificar tachando
\usepackage{polynom}    % Divisiones y Ruffini
\usepackage{units} % Para poner fracciones diagonales con \nicefrac

\usepackage{pgfplots}   %Representar funciones
\pgfplotsset{compat=1.18}  % Versión 1.18

\usepackage{tikz-cd}    % Para usar diagramas de composiciones
\usetikzlibrary{calc}   % Para usar cálculo de coordenadas en tikz

%Definición de teoremas, etc.
\usepackage{amsthm}
%\swapnumbers   % Intercambia la posición del texto y de la numeración

\theoremstyle{plain}

\makeatletter
\@ifclassloaded{article}{
  \newtheorem{teo}{Teorema}[section]
}{
  \newtheorem{teo}{Teorema}[chapter]  % Se resetea en cada chapter
}
\makeatother

\newtheorem{coro}{Corolario}[teo]           % Se resetea en cada teorema
\newtheorem{prop}[teo]{Proposición}         % Usa el mismo contador que teorema
\newtheorem{lema}[teo]{Lema}                % Usa el mismo contador que teorema

\theoremstyle{remark}
\newtheorem*{observacion}{Observación}

\theoremstyle{definition}

\makeatletter
\@ifclassloaded{article}{
  \newtheorem{definicion}{Definición} [section]     % Se resetea en cada chapter
}{
  \newtheorem{definicion}{Definición} [chapter]     % Se resetea en cada chapter
}
\makeatother

\newtheorem*{notacion}{Notación}
\newtheorem*{ejemplo}{Ejemplo}
\newtheorem*{ejercicio*}{Ejercicio}             % No numerado
\newtheorem{ejercicio}{Ejercicio} [section]     % Se resetea en cada section


% Modificar el formato de la numeración del teorema "ejercicio"
\renewcommand{\theejercicio}{%
  \ifnum\value{section}=0 % Si no se ha iniciado ninguna sección
    \arabic{ejercicio}% Solo mostrar el número de ejercicio
  \else
    \thesection.\arabic{ejercicio}% Mostrar número de sección y número de ejercicio
  \fi
}


% \renewcommand\qedsymbol{$\blacksquare$}         % Cambiar símbolo QED
%------------------------------------------------------------------------

% Paquetes para encabezados
\usepackage{fancyhdr}
\pagestyle{fancy}
\fancyhf{}

\newcommand{\helv}{ % Modificación tamaño de letra
\fontfamily{}\fontsize{12}{12}\selectfont}
\setlength{\headheight}{15pt} % Amplía el tamaño del índice


%\usepackage{lastpage}   % Referenciar última pag   \pageref{LastPage}
\fancyfoot[C]{\thepage}

%------------------------------------------------------------------------

% Conseguir que no ponga "Capítulo 1". Sino solo "1."
\makeatletter
\@ifclassloaded{book}{
  \renewcommand{\chaptermark}[1]{\markboth{\thechapter.\ #1}{}} % En el encabezado
    
  \renewcommand{\@makechapterhead}[1]{%
  \vspace*{50\p@}%
  {\parindent \z@ \raggedright \normalfont
    \ifnum \c@secnumdepth >\m@ne
      \huge\bfseries \thechapter.\hspace{1em}\ignorespaces
    \fi
    \interlinepenalty\@M
    \Huge \bfseries #1\par\nobreak
    \vskip 40\p@
  }}
}
\makeatother

%------------------------------------------------------------------------
% Paquetes de cógido
\usepackage{minted}
\renewcommand\listingscaption{Código fuente}

\usepackage{fancyvrb}
% Personaliza el tamaño de los números de línea
\renewcommand{\theFancyVerbLine}{\small\arabic{FancyVerbLine}}

% Estilo para C++
\newminted{cpp}{
    frame=lines,
    framesep=2mm,
    baselinestretch=1.2,
    linenos,
    escapeinside=||
}



\usepackage{listings} % Para incluir código desde un archivo

\renewcommand\lstlistingname{Código Fuente}
\renewcommand\lstlistlistingname{Índice de Códigos Fuente}

% Definir colores
\definecolor{vscodepurple}{rgb}{0.5,0,0.5}
\definecolor{vscodeblue}{rgb}{0,0,0.8}
\definecolor{vscodegreen}{rgb}{0,0.5,0}
\definecolor{vscodegray}{rgb}{0.5,0.5,0.5}
\definecolor{vscodebackground}{rgb}{0.97,0.97,0.97}
\definecolor{vscodelightgray}{rgb}{0.9,0.9,0.9}

% Configuración para el estilo de C similar a VSCode
\lstdefinestyle{vscode_C}{
  backgroundcolor=\color{vscodebackground},
  commentstyle=\color{vscodegreen},
  keywordstyle=\color{vscodeblue},
  numberstyle=\tiny\color{vscodegray},
  stringstyle=\color{vscodepurple},
  basicstyle=\scriptsize\ttfamily,
  breakatwhitespace=false,
  breaklines=true,
  captionpos=b,
  keepspaces=true,
  numbers=left,
  numbersep=5pt,
  showspaces=false,
  showstringspaces=false,
  showtabs=false,
  tabsize=2,
  frame=tb,
  framerule=0pt,
  aboveskip=10pt,
  belowskip=10pt,
  xleftmargin=10pt,
  xrightmargin=10pt,
  framexleftmargin=10pt,
  framexrightmargin=10pt,
  framesep=0pt,
  rulecolor=\color{vscodelightgray},
  backgroundcolor=\color{vscodebackground},
}

%------------------------------------------------------------------------

% Comandos definidos
\newcommand{\bb}[1]{\mathbb{#1}}
\newcommand{\cc}[1]{\mathcal{#1}}

% I prefer the slanted \leq
\let\oldleq\leq % save them in case they're every wanted
\let\oldgeq\geq
\renewcommand{\leq}{\leqslant}
\renewcommand{\geq}{\geqslant}

% Si y solo si
\newcommand{\sii}{\iff}

% Letras griegas
\newcommand{\eps}{\epsilon}
\newcommand{\veps}{\varepsilon}
\newcommand{\lm}{\lambda}

\newcommand{\ol}{\overline}
\newcommand{\ul}{\underline}
\newcommand{\wt}{\widetilde}
\newcommand{\wh}{\widehat}

\let\oldvec\vec
\renewcommand{\vec}{\overrightarrow}

% Derivadas parciales
\newcommand{\del}[2]{\frac{\partial #1}{\partial #2}}
\newcommand{\Del}[3]{\frac{\partial^{#1} #2}{\partial^{#1} #3}}
\newcommand{\deld}[2]{\dfrac{\partial #1}{\partial #2}}
\newcommand{\Deld}[3]{\dfrac{\partial^{#1} #2}{\partial^{#1} #3}}


\newcommand{\AstIg}{\stackrel{(\ast)}{=}}
\newcommand{\Hop}{\stackrel{L'H\hat{o}pital}{=}}

\newcommand{\red}[1]{{\color{red}#1}} % Para integrales, destacar los cambios.

% Método de integración
\newcommand{\MetInt}[2]{
    \left[\begin{array}{c}
        #1 \\ #2
    \end{array}\right]
}

% Declarar aplicaciones
% 1. Nombre aplicación
% 2. Dominio
% 3. Codominio
% 4. Variable
% 5. Imagen de la variable
\newcommand{\Func}[5]{
    \begin{equation*}
        \begin{array}{rrll}
            #1:& #2 & \longrightarrow & #3\\
               & #4 & \longmapsto & #5
        \end{array}
    \end{equation*}
}

%------------------------------------------------------------------------
\let\oldRe\Re % save them in case they're every wanted
\let\oldIm\Im
\renewcommand{\Re}{\operatorname{Re}} % redefine them
\renewcommand{\Im}{\operatorname{Im}}
\DeclareMathOperator{\Log}{Log}
\DeclareMathOperator{\Arg}{Arg}

\begin{document}

    % 1. Foto de fondo
    % 2. Título
    % 3. Encabezado Izquierdo
    % 4. Color de fondo
    % 5. Coord x del titulo
    % 6. Coord y del titulo
    % 7. Fecha

    
    % 1. Foto de fondo
% 2. Título
% 3. Encabezado Izquierdo
% 4. Color de fondo
% 5. Coord x del titulo
% 6. Coord y del titulo
% 7. Fecha

\newcommand{\portada}[7]{

    \portadaBase{#1}{#2}{#3}{#4}{#5}{#6}{#7}
    \portadaBook{#1}{#2}{#3}{#4}{#5}{#6}{#7}
}

\newcommand{\portadaExamen}[7]{

    \portadaBase{#1}{#2}{#3}{#4}{#5}{#6}{#7}
    \portadaArticle{#1}{#2}{#3}{#4}{#5}{#6}{#7}
}




\newcommand{\portadaBase}[7]{

    % Tiene la portada principal y la licencia Creative Commons
    
    % 1. Foto de fondo
    % 2. Título
    % 3. Encabezado Izquierdo
    % 4. Color de fondo
    % 5. Coord x del titulo
    % 6. Coord y del titulo
    % 7. Fecha
    
    
    \thispagestyle{empty}               % Sin encabezado ni pie de página
    \newgeometry{margin=0cm}        % Márgenes nulos para la primera página
    
    
    % Encabezado
    \fancyhead[L]{\helv #3}
    \fancyhead[R]{\helv \nouppercase{\leftmark}}
    
    
    \pagecolor{#4}        % Color de fondo para la portada
    
    \begin{figure}[p]
        \centering
        \transparent{0.3}           % Opacidad del 30% para la imagen
        
        \includegraphics[width=\paperwidth, keepaspectratio]{assets/#1}
    
        \begin{tikzpicture}[remember picture, overlay]
            \node[anchor=north west, text=white, opacity=1, font=\fontsize{60}{90}\selectfont\bfseries\sffamily, align=left] at (#5, #6) {#2};
            
            \node[anchor=south east, text=white, opacity=1, font=\fontsize{12}{18}\selectfont\sffamily, align=right] at (9.7, 3) {\textbf{\href{https://losdeldgiim.github.io/}{Los Del DGIIM}}};
            
            \node[anchor=south east, text=white, opacity=1, font=\fontsize{12}{15}\selectfont\sffamily, align=right] at (9.7, 1.8) {Doble Grado en Ingeniería Informática y Matemáticas\\Universidad de Granada};
        \end{tikzpicture}
    \end{figure}
    
    
    \restoregeometry        % Restaurar márgenes normales para las páginas subsiguientes
    \pagecolor{white}       % Restaurar el color de página
    
    
    \newpage
    \thispagestyle{empty}               % Sin encabezado ni pie de página
    \begin{tikzpicture}[remember picture, overlay]
        \node[anchor=south west, inner sep=3cm] at (current page.south west) {
            \begin{minipage}{0.5\paperwidth}
                \href{https://creativecommons.org/licenses/by-nc-nd/4.0/}{
                    \includegraphics[height=2cm]{assets/Licencia.png}
                }\vspace{1cm}\\
                Esta obra está bajo una
                \href{https://creativecommons.org/licenses/by-nc-nd/4.0/}{
                    Licencia Creative Commons Atribución-NoComercial-SinDerivadas 4.0 Internacional (CC BY-NC-ND 4.0).
                }\\
    
                Eres libre de compartir y redistribuir el contenido de esta obra en cualquier medio o formato, siempre y cuando des el crédito adecuado a los autores originales y no persigas fines comerciales. 
            \end{minipage}
        };
    \end{tikzpicture}
    
    
    
    % 1. Foto de fondo
    % 2. Título
    % 3. Encabezado Izquierdo
    % 4. Color de fondo
    % 5. Coord x del titulo
    % 6. Coord y del titulo
    % 7. Fecha


}


\newcommand{\portadaBook}[7]{

    % 1. Foto de fondo
    % 2. Título
    % 3. Encabezado Izquierdo
    % 4. Color de fondo
    % 5. Coord x del titulo
    % 6. Coord y del titulo
    % 7. Fecha

    % Personaliza el formato del título
    \pretitle{\begin{center}\bfseries\fontsize{42}{56}\selectfont}
    \posttitle{\par\end{center}\vspace{2em}}
    
    % Personaliza el formato del autor
    \preauthor{\begin{center}\Large}
    \postauthor{\par\end{center}\vfill}
    
    % Personaliza el formato de la fecha
    \predate{\begin{center}\huge}
    \postdate{\par\end{center}\vspace{2em}}
    
    \title{#2}
    \author{\href{https://losdeldgiim.github.io/}{Los Del DGIIM}}
    \date{Granada, #7}
    \maketitle
    
    \tableofcontents
}




\newcommand{\portadaArticle}[7]{

    % 1. Foto de fondo
    % 2. Título
    % 3. Encabezado Izquierdo
    % 4. Color de fondo
    % 5. Coord x del titulo
    % 6. Coord y del titulo
    % 7. Fecha

    % Personaliza el formato del título
    \pretitle{\begin{center}\bfseries\fontsize{42}{56}\selectfont}
    \posttitle{\par\end{center}\vspace{2em}}
    
    % Personaliza el formato del autor
    \preauthor{\begin{center}\Large}
    \postauthor{\par\end{center}\vspace{3em}}
    
    % Personaliza el formato de la fecha
    \predate{\begin{center}\huge}
    \postdate{\par\end{center}\vspace{5em}}
    
    \title{#2}
    \author{\href{https://losdeldgiim.github.io/}{Los Del DGIIM}}
    \date{Granada, #7}
    \thispagestyle{empty}               % Sin encabezado ni pie de página
    \maketitle
    \vfill
}
    \portadaExamen{ffccA4.jpg}{Variable Compleja I\\Examen III}{Variable Compleja I. Examen III}{MidnightBlue}{-8}{28}{2024-2025}{Arturo Olivares Martos}

    \begin{description}
        \item[Asignatura] Variable Compleja I.
        \item[Curso Académico] 2022-23.
        \item[Grado] Doble Grado en Ingeniería Informática y Matemáticas.
        \item[Grupo] Único.
        \item[Profesor] Javier Merí de la Maza.
        \item[Descripción] Prueba Intermedia.
        \item[Fecha] 9 de Mayo de 2023.
        \item[Duración] 120 minutos.
    
    \end{description}
    \newpage

    \begin{ejercicio}[4 puntos]
        Probar que que la serie $\displaystyle \sum_{n\geq 0} \frac{\sen(nz)}{2^n}$ converge absolutamente en todo punto del dominio $\Omega=\{z \in \mathbb{C} : -\ln(2) < \Im z < \ln(2)\}$. Estudiar la convergencia uniforme en subconjuntos de $\Omega$. Probar que la función $g$ siguiente es continua, y calcular la integral $\displaystyle \int_{C(0,1)} \frac{g(z)}{z}\ dz$:
        \Func{g}{\Omega}{\bb{C}}{z}{\displaystyle \sum\limits_{n=0}^{\infty} \frac{\sen(nz)}{2^n}}
    \end{ejercicio}

    \begin{ejercicio}[2 puntos]
        Estudiar la derivabilidad de la función $f$ siguiente:
        \Func{f}{\bb{C}}{\bb{C}}{z}{e^{\ol{z}}}
    \end{ejercicio}

    \begin{ejercicio}
        Sea $\Omega\subseteq \mathbb{R}^2 \equiv\mathbb{C}$ un abierto. Decimos que una función $\varphi:\Omega\to\mathbb{R}$ es armónica en $\Omega$ si $\varphi \in C^2(\Omega)$ y
        \begin{equation*}
            \frac{\partial^2 \varphi}{\partial x^2}(x,y)+\frac{\partial^2 \varphi}{\partial y^2}(x,y) = 0 \quad \forall (x,y) \in \Omega.
        \end{equation*}
        \begin{enumerate}
            \item \textbf{[1.5 puntos]} Sea $\varphi:\Omega\to\mathbb{R}$ armónica en $\Omega$. Probar que la función $g$ siguiente es holomorfa en $\Omega$:
            \Func{g}{\Omega}{\bb{C}}{x+iy}{\frac{\partial \varphi}{\partial x}(x,y)-i\frac{\partial \varphi}{\partial y}(x,y)}
            \item \textbf{[1.5 puntos]} Suponiendo que $\Omega$ es un dominio estrellado, probar que existe $f \in \cc{H}(\Omega)$ de modo que $\Re f = \varphi$ y que $f$ es única salvo una constante.
            \item \textbf{[1 punto]} Deducir que una función armónica en un dominio estrellado es, de hecho, de clase $C^{\infty}$.
        \end{enumerate}
    \end{ejercicio}

    \newpage
    \setcounter{ejercicio}{0}

    \begin{ejercicio}[4 puntos]
        Probar que que la serie $\displaystyle \sum_{n\geq 0} \frac{\sen(nz)}{2^n}$ converge absolutamente en todo punto del dominio $\Omega=\{z \in \mathbb{C} : -\ln(2) < \Im z < \ln(2)\}$. Estudiar la convergencia uniforme en subconjuntos de $\Omega$. Probar que la función $g$ siguiente es continua, y calcular la integral $\displaystyle \int_{C(0,1)} \frac{g(z)}{z}\ dz$:
        \Func{g}{\Omega}{\bb{C}}{z}{\displaystyle \sum\limits_{n=0}^{\infty} \frac{\sen(nz)}{2^n}}

        Probemos en primer lugar que, dado $A\subset \Omega$, la serie converge uniformemente en $A$ si y solo si:
        \begin{equation*}
            \sup\{|\Im z| : z\in A\} < \ln(2)
        \end{equation*}
        \begin{description}
            \item[$\Longrightarrow)$] Demostraremos el recíproco. Como $A\subset \Omega$, no puede darse que dicho supremo sea mayor que $\ln(2)$. Supongamos por tanto que dicho supremo es $2$. Entonces, por la caracterización del supremo mediante sucesiones, existe una sucesión $\left\{\wt{z}_n\right\}$ de puntos de $A$ tal que $\left\{|\Im \wt{z}_n|\right\}\to \ln 2$. Tomando una parcial suya, existe una sucesión $\left\{z_n\right\}$ de puntos de $A$ tal que $\left|\Im z_n\right|>\ln 2 - \nicefrac{1}{n}$ para todo $n\in \bb{N}$.\\
            
            Para cada $n\in \bb{N}$, tenemos que:
            \begin{align*}
                \left|\dfrac{\sen(nz_n)}{2^n}\right| &\geq \dfrac{\left|\senh(n\Im z_n)\right|}{2^{n}}
                = \dfrac{\senh(n|\Im z_n|)}{2^{n}}
                = \dfrac{e^{n|\Im z_n|}-e^{-n|\Im z_n|}}{2^{n+1}} \geq \dfrac{e^{n|\Im z_n|}}{2^{n+1}}
                > \\
                &> \dfrac{e^{n\ln 2 -1}}{2^{n+1}}
                = \dfrac{1}{2e}
            \end{align*}

            Por tanto, como hemos encontrado una sucesión de puntos de $A$ tal que el término general no converge a $0$, entonces el término general no converge uniformemente a $0$. Por tanto, la serie no converge uniformemente en $A$.

            \item[$\Longleftarrow)$] Sea dicho supremo $r<\ln 2$. Entonces, para todo $z\in A$ y todo $n\in \bb{N}$, tenemos que:
            \begin{align*}
                \left|\dfrac{\sen(nz)}{2^n}\right| &<\left(\dfrac{e^{|\Im z|}}{2}\right)^n
                < \left(\dfrac{e^{r}}{2}\right)^n
            \end{align*}

            Como $r<\ln(2)$, tenemos que $e^{r}<2$, por lo que la base de la potencia es menor que 1. Por tanto, la serie geométrica de razón dicho cociente es convergente. Por el Test de Weierstrass, la serie de partida converge uniformemente en $A$.
        \end{description}

        En particular, dado $K\subset \Omega$ compacto, en particular es cerrado, y el supremo en cuestión se alcanza. Como es un subconjunto de $\Omega$, entonces:
        \begin{align*}
            \sup\{|\Im w| : w\in K\}=\max\{|\Im w| : w\in K\} < \ln(2)
        \end{align*}
        Por tanto, la serie converge uniformemente en compactos de $\Omega$.\\

        Para la convergencia absoluta, consideramos $z\in \Omega$ fijo. Como $\{z\}$ es compacto, por el Test de Weierstrass tenemos que converge absolutamente en $\{z\}$. Como $z$ es arbitrario, tenemos que la serie converge absolutamente en todo punto de $\Omega$.\\

        Para probar que $g$ es continua, consideramos $z\in \Omega$. Como $\Omega$ es abierto, existe $R\in \bb{R}^+$ tal que $D(z,R)\subset \Omega$. Considerando ahora $\nicefrac{R}{2}$, tenemos que $z\in \ol{D}(z,\nicefrac{R}{2})\subset \Omega$. Por ser este compacto, la serie converge uniformemente en $\ol{D}(z,\nicefrac{R}{2})$. Como el término general de la serie es continuo para todo $n\in \bb{N}$, entonces $g$ es continua en $z$. Como $z$ es arbitrario, $g$ es continua en $\Omega$.\\

        Estudiemos ahora la convergencia uniforme del integrando en $C(0,1)$. Como $C(0,1)$ es compacto, entonces $g$ converge uniformemente y $z\mapsto \nicefrac{1}{z}$ es acotado. Por tanto, el integrando converge uniformemente en $C(0,1)$, luego podemos intercambiar la sumatoria y la integral, algo que nos será de utilidad a la hora de calcular la integral:
        \begin{align*}
            \int_{C(0,1)} \frac{g(z)}{z}\ dz &= \sum_{n=0}^{\infty} \int_{C(0,1)} \frac{\sen(nz)}{z2^n}\ dz
            = \sum_{n=0}^{\infty} \frac{1}{2^n}\int_{C(0,1)} \frac{\sen(nz)}{z}\ dz \AstIg\\&\AstIg \sum_{n=0}^{\infty} \frac{1}{2^n}\cdot 2\pi i \cdot \sen (n\cdot 0) = \sum_{n=0}^{\infty} 0 = 0
        \end{align*}
        donde en $(\ast)$ hemos aplicado la Fórmula de Cauchy para la circunferencia, usando que el seno es una función entera.
    \end{ejercicio}

    \begin{ejercicio}[2 puntos]
        Estudiar la derivabilidad de la función $f$ siguiente:
        \Func{g}{\bb{C}}{\bb{C}}{z}{e^{\ol{z}}}


        Sea $z=x+iy\in \bb{C}$. Definimos $u,v:\mathbb{R}^2\to\mathbb{R}$ como:
        \begin{align*}
            u(x,y) &= \Re(g(x+iy)) = \Re(e^{x-iy}) = e^x\cos (-y)\\
            v(x,y) &= \Im(g(x+iy)) = \Im(e^{x-iy}) = e^x\sen (-y)
        \end{align*}

        Calculamos ahora las derivadas parciales de $u$ y $v$:
        \begin{align*}
            \frac{\partial u}{\partial x}(x,y) = e^x\cos (-y)&
            \frac{\partial u}{\partial y}(x,y) = e^x\sen (-y)\\
            \frac{\partial v}{\partial x}(x,y) = e^x\sen (-y)&
            \frac{\partial v}{\partial y}(x,y) = -e^x\cos (-y)
        \end{align*}

        La primera ecuación de Cauchy-Riemann nos dice:
        \begin{align*}
            \frac{\partial u}{\partial x}(x,y) = e^x\cos(-y)
            = \frac{\partial v}{\partial y}(x,y) = -e^x\cos(-y)
        \end{align*}

        Para que esto se de, es necesario que $e^x=0$ (lo cual no se da) o que $\cos(-y)=cos(y)=0$.\\

        La segunda ecuación de Cauchy-Riemann nos dice:
        \begin{align*}
            \frac{\partial u}{\partial y}(x,y) = e^x\sen(-y)
            = -\frac{\partial v}{\partial x}(x,y) = -e^x\sen(-y)
        \end{align*}
        Para que esto se de, es necesario que $e^x=0$ (lo cual no se da) o que $\sen(-y)=-\sen(y)=0$.

        Como no es posible que el seno y el coseno reales se anulen simultáneamente, se concluye que $g$ no es derivable en $z$ (puesto que no se cumplen las Ecuaciones de Cauchy-Riemann).
    \end{ejercicio}

    \begin{ejercicio}
        Sea $\Omega\subseteq \mathbb{R}^2 \equiv\mathbb{C}$ un abierto. Decimos que una función $\varphi:\Omega\to\mathbb{R}$ es armónica en $\Omega$ si $\varphi \in C^2(\Omega)$ y
        \begin{equation*}
            \dfrac{\partial^2 \varphi}{\partial x^2}(x,y)+\dfrac{\partial^2 \varphi}{\partial y^2}(x,y) = 0 \quad \forall (x,y) \in \Omega.
        \end{equation*}
        \begin{enumerate}
            \item \textbf{[1.5 puntos]} Sea $\varphi:\Omega\to\mathbb{R}$ armónica en $\Omega$. Probar que la función $g$ siguiente es holomorfa en $\Omega$:
            \Func{g}{\Omega}{\bb{C}}{x+iy}{\frac{\partial \varphi}{\partial x}(x,y)-i\frac{\partial \varphi}{\partial y}(x,y)}

            Definimos las funciones $u,v:\Omega\to \bb{R}$, con vistas a aplicar las Ecuaciones de Cauchy-Riemann:
        \begin{align*}
            u(x,y) &= \Re g(x+iy) = \frac{\partial \varphi}{\partial x}(x, y)\\
            v(x,y) &= \Im g(x+iy) = -\frac{\partial \varphi}{\partial y}(x, y)
        \end{align*}

        Calculamos cada una de las derivadas parciales de $u$ y $v$:
        \begin{align*}
            \frac{\partial u}{\partial x} &= \frac{\partial^2 \varphi}{\partial x^2}&
            \frac{\partial u}{\partial y} = \frac{\partial^2 \varphi}{\partial x \partial y}\\
            \frac{\partial v}{\partial x} &= -\frac{\partial^2 \varphi}{\partial y \partial x}&
            \frac{\partial v}{\partial y} = -\frac{\partial^2 \varphi}{\partial y^2}
        \end{align*}

        Comprobemos que se cumplen las Ecuaciones de Cauchy-Riemann:
        \begin{align*}
            \frac{\partial u}{\partial x} &= \frac{\partial^2 \varphi}{\partial x^2} = -\frac{\partial^2 \varphi}{\partial y^2} = \frac{\partial v}{\partial y}\iff \frac{\partial^2 \varphi}{\partial x^2} + \frac{\partial^2 \varphi}{\partial y^2} = 0\\
            \frac{\partial u}{\partial y} &= \frac{\partial^2 \varphi}{\partial x \partial y} = -\left(-\frac{\partial^2 \varphi}{\partial y \partial x}\right) = -\frac{\partial v}{\partial x}
        \end{align*}
        Notemos que la primera ecuación se cumple por hipótesis, y la segunda se cumple por el Teorema de Clairaut. Por tanto, se cumplen las Ecuaciones de Cauchy-Riemann, y por tanto $g$ es holomorfa en $\Omega$.

            \item \textbf{[1.5 puntos]} Suponiendo que $\Omega$ es un dominio estrellado, probar que existe $f \in \cc{H}(\Omega)$ de modo que $\Re f = \varphi$ y que $f$ es única salvo una constante.
            
            Como $\Omega$ es un dominio estrellado y $g\in \cc{H}(\Omega)$, por el Teorema Local de Cauchy, $\exists f\in \cc{H}(\Omega)$ una primitiva de $g$ en $\Omega$. Por tanto, tenemos que:
        \begin{align*}
            f'(x+iy) &= g(x + iy) = \frac{\partial \varphi}{\partial x}(x, y) - i\cdot  \frac{\partial \varphi}{\partial y}(x, y)\qquad \forall (x, y) \in \Omega
        \end{align*}

        Definimos ahora de nuevo $u,v:\Omega\to \bb{R}$ como:
        \begin{align*}
            u(x,y) &= \Re f(x+iy)\\
            v(x,y) &= \Im f(x+iy)
        \end{align*}

        Por las Ecuaciones de Cauchy-Riemann, tenemos que:
        \begin{align*}
            \frac{\partial u}{\partial x} &= \frac{\partial v}{\partial y}= \frac{\partial \varphi}{\partial x}\\
            \frac{\partial v}{\partial x} &= -\frac{\partial \varphi}{\partial y}=-\frac{\partial u}{\partial y}
        \end{align*}

        Para calcular $u=\Re f$, integramos la primera ecuación respecto de $x$:
        \begin{align*}
            u(x,y) &= \int \frac{\partial \varphi}{\partial x}\ dx = \varphi(x,y) + h(y)\qquad \forall (x,y) \in \Omega
        \end{align*}
        donde $h$ es una función de $y$ que no depende de $x$ y representa la constante de integración. Derivando respecto de $y$, tenemos que:
        \begin{align*}
            \frac{\partial u}{\partial y}(x,y) &= \frac{\partial \varphi}{\partial y}(x,y) + h'(y)\Longrightarrow
            h'(y) = 0\qquad \forall (x,y) \in \Omega
        \end{align*}

        Por tanto, $h$ es constante, por lo que $\exists C\in \bb{R}$ tal que:
        \begin{align*}
            u(x,y) &= \varphi(x,y) + C\qquad \forall (x,y) \in \Omega
        \end{align*}  
        
        Como $u=\Re f$, y por hipótesis buscamos que $\Re f = \varphi$, tenemos que $C=0$. Por tanto, tenemos que:
        \begin{align*}
            \Re f(x,y) &= \varphi(x,y)\qquad \forall (x,y) \in \Omega
        \end{align*}


        Por tanto, la existencia está probada. Supongamos ahora que existe otra función $g\in \cc{H}(\Omega)$ tal que $\Re g = \varphi$. Definimos $h = f - g\in \cc{H}(\Omega)$. Entonces, tenemos que:
        \begin{align*}
            \Re h &= \Re f - \Re g = \varphi - \varphi = 0
        \end{align*}

        Como $h$ es holomorfa, está definida en un dominio, y su parte real es nula, entonces $h$ es constante, luego $\exists \lambda \in \bb{C}$ tal que:
        \begin{align*}
            f(z) &= g(z) + \lambda\qquad \forall z \in \Omega
        \end{align*}

        Por tanto, $f$ es única salvo una constante.

            \item \textbf{[1 punto]} Deducir que una función armónica en un dominio estrellado es, de hecho, de clase $C^{\infty}$.\\
            
            Por el apartado anterior, sabemos que $\exists f\in \cc{H}(\Omega)$ tal que $\Re f = \varphi$. Como $f$ es holomorfa, es infinitamente derivable. Calculemos su primera derivada empleando las Ecuaciones de Cauchy-Riemann y que $\Re f = \varphi$:
            \begin{align*}
                f'(x+iy) &= \frac{\partial \varphi}{\partial x}(x,y) - i\cdot \frac{\partial \varphi}{\partial y}(x,y)\qquad \forall (x,y) \in \Omega
            \end{align*}

            Como $f'\in \cc{H}(\Omega)$, en particular es continua, por lo que las derivadas de orden $1$ de $\varphi$ son continuas. Por tanto, $\varphi\in C^1(\Omega)$.\\

            Aplicamos ahora las ecuaciones de Cauchy-Riemann a la función $f'$:
            \begin{align*}
                f''(x+iy) &= \frac{\partial^2 \varphi}{\partial x^2}(x,y) - i\cdot \frac{\partial^2 \varphi}{\partial x \partial y}(x,y)\\
                &= -\frac{\partial^2 \varphi}{\partial y^2}(x,y) - i\cdot \frac{\partial^2 \varphi}{\partial y \partial x}(x,y)\qquad \forall (x,y) \in \Omega
            \end{align*}

            Como $f''\in \cc{H}(\Omega)$, en particular es continua, por lo que las $4$ derivadas de orden $2$ de $\varphi$ son continuas. Por tanto, $\varphi\in C^2(\Omega)$.\\

            Aplicando de nuevo las Ecuaciones de Cauchy-Riemann a la función $f''$ para obtener $f'''$, podríamos obtener que las $8$ derivadas de orden $3$ de $\varphi$ son continuas. Por tanto, $\varphi\in C^3(\Omega)$; y así sucesivamente.
        \end{enumerate}
    \end{ejercicio}

\end{document}